% ==============================================================================
% Section 3: Theoretical Framework
% Paper: Inside Synthetic Dollar Liquidity and the Microstructure of FX Swap Markets
% ==============================================================================

\section{Theoretical Framework}
\label{sec:theory}

Dollar funding crises are not runs against the dollar. They are runs \emph{from}
inside synthetic dollar liquidity into outside safe dollar liquidity. This section
formalizes that claim. We construct an equilibrium model of the FX swap market
in which the conversion premium between synthetic and safe dollars is determined
by a matching process between dollar seekers and conversion intermediaries. Three
results follow: demand pressure raises the equilibrium premium (Proposition~1),
capacity tightening amplifies the effect (Proposition~2), and the premium
responds convexly to market tightness under plausible conditions on the quality
function (Proposition~3). The first two propositions map directly to the IV and
interaction specifications in \cref{sec:results}; the third characterizes the
theoretical shape of price impact and motivates our descriptive evidence on
crisis episodes.


% ------------------------------------------------------------------------------
\subsection{Environment: Inside and Outside Dollar Liquidity}
\label{subsec:environment}
% ------------------------------------------------------------------------------

We model two forms of dollar liquidity that differ in their institutional
connection to the US monetary core. \emph{Outside safe dollar liquidity}, $O_t$,
encompasses Federal Reserve reserve balances, highly liquid Treasury securities,
and domestic bank liabilities backstopped by FDIC insurance or discount window
access, as well as dollar liquidity extended through central-bank swap lines.
\emph{Inside synthetic dollar liquidity}, $I_t$, is dollar liquidity produced
outside the US domestic banking system---primarily through FX swaps, offshore
repo, Eurodollar deposits, and the wholesale funding operations of global banks.
A unit of $I_t$ is obtained when a non-US entity delivers domestic currency on
the near leg of an FX swap and commits to return dollars on the far leg.

The distinction matters because in normal times $I_t$ and $O_t$ are near-perfect
substitutes: arbitrage keeps the cross-currency basis close to zero. In stress,
they are not. When the probability of finding a counterparty able to supply
outside dollars at short notice falls, the effective service that one unit of
$I_t$ delivers deteriorates. We capture this through an endogenous quality
parameter $\chi_t \in (0,1]$, defined as the number of units of outside safe
dollar service that one unit of $I_t$ delivers given the current state of the
conversion market. When every seeker finds a counterparty instantly, $\chi_t = 1$
and inside liquidity is a perfect substitute. When matching is difficult,
$\chi_t < 1$ and agents pay a premium to convert.

The aggregate dollar liquidity service available to agents is produced by a CES
technology:
\begin{equation}
  L_t = \bigl[\alpha_t O_t^{\rho} + (1-\alpha_t)(\chi_t I_t)^{\rho}
  \bigr]^{1/\rho},
  \label{eq:CES}
\end{equation}
where $\sigma \equiv 1/(1-\rho) > 1$ is the elasticity of substitution.
The parameter $\sigma > 1$ captures the gross-substitutes property:
a higher premium on inside dollars induces substitution toward outside
dollars. Under $\sigma > 1$, inside and outside dollars are not complements,
and a stress shock that degrades $\chi_t$ triggers demand reallocation---the
run-from-inside mechanism. The condition $\sigma > 1$ is maintained as a
testable hypothesis: if the demand system estimation yields $\mathcal{B} =
\sum_s \beta_s \geq 0$ (upward-sloping aggregate net demand), the substitution
assumption is rejected by the data.%
\footnote{This test is reported in \cref{subsec:demand-system}. Failure would
imply that inside and outside dollar liquidity are gross complements in our
sample, which would undermine the run-from-inside interpretation.}

The FX swap market is where the conversion of $I_t$ to $O_t$ is priced and
intermediated. Three institutional actors structure this market. \emph{Foreign
banks} (and their FBO branches in the US) are the primary dollar seekers: they
raise foreign-currency funding and convert it to dollars synthetically to meet
dollar payment obligations. \emph{Primary dealers, US commercial banks, and FBO
branches} are the primary converters: they supply dollars against foreign
currency, intermediating between the dollar money market (outside liquidity) and
foreign-currency-funded balance sheets. \emph{Hedge funds} enter as a secondary
supplier when the basis is wide enough to make the basis trade profitable; they
withdraw in acute stress when risk capacity is exhausted.

The equilibrium funding spread $\mu_t \equiv R_{I,t} - R_{O,t}$, where $R_{I,t}$
and $R_{O,t}$ are the gross costs of inside and outside dollar liquidity, is the
primary object of interest. In the data, $\mu_{m,t} = \text{PriceUSD}_{m,t} =
-\text{TIB}_{m,t}$, where TIB is the transaction-implied basis for currency pair
$m$. A positive spread means inside dollar liquidity is more expensive than
outside: agents pay a premium for synthetic dollars.


% ------------------------------------------------------------------------------
\subsection{Matching Technology}
\label{subsec:matching}
% ------------------------------------------------------------------------------

Agents who wish to convert inside to outside dollar liquidity must find a
counterparty willing to supply outside dollars. This is a search process: not
every seeker finds a counterparty in every period. We model the conversion market
as a standard search-and-matching market.

Let $U_t \geq 0$ denote the mass of agents actively seeking conversion in period
$t$---the demand side---and $B_t \geq 0$ the aggregate balance-sheet capacity
of conversion intermediaries---the supply side. Total matches are given by the
Cobb-Douglas matching function
\begin{equation}
  M_t = \eta\, U_t^{\gamma}\, B_t^{1-\gamma},
  \label{eq:matching}
\end{equation}
where $\eta > 0$ is matching efficiency and $\gamma \in (0,1)$ is the
demand-side elasticity of matching. The function satisfies constant returns to
scale in $(U_t, B_t)$ and is increasing and concave in each argument. The
matching probability for a given seeker is
\begin{equation}
  p_t = \frac{M_t}{U_t} = \eta \left(\frac{B_t}{U_t}\right)^{1-\gamma}
  = \eta\,\theta_t^{-(1-\gamma)},
  \label{eq:match-prob}
\end{equation}
where $\theta_t \equiv U_t / B_t$ is \emph{market tightness}. When $\theta_t$
is high---many seekers relative to intermediary capacity---each seeker is less
likely to find a counterparty. This is the search externality: each additional
seeker reduces the matching probability for all others.

The matching probability determines quality: $\chi_t = \chi(p_t)$, where $\chi$
is strictly increasing with $\chi(1) = 1$ and $\chi'(p) > 0$. A lower
matching probability thus degrades the effective service that inside dollar
liquidity provides.

\begin{lemma}[Matching Function Properties]
  \label{lem:matching}
  The matching function $M_t = \eta U_t^{\gamma} B_t^{1-\gamma}$ satisfies:
  (i) $M_t \leq \min(U_t, B_t)$ for $\eta \leq 1$ (feasibility); (ii) $M_t$
  is increasing in both $U_t$ and $B_t$; (iii) $p_t = M_t / U_t$ is strictly
  decreasing in $\theta_t = U_t / B_t$; (iv) $\partial p_t / \partial B_t > 0$:
  additional intermediary capacity raises the matching probability for seekers.
\end{lemma}

Properties (iii) and (iv) are the model's core mechanism. A rise in $U_t$ (more
seekers) tightens the market and reduces $p_t$; a rise in $B_t$ (more capacity)
loosens it and raises $p_t$. Both work through $\theta_t$.

Intermediary capacity is the sum of contributions from five institutional types:
\begin{equation}
  B_t = B_t^{FBO} + B_t^{USbank} + B_t^{dealer} + B_t^{HF} + B_t^{CBswap}.
  \label{eq:capacity}
\end{equation}
FBO branches and primary dealers supply capacity that is (weakly) decreasing in
financial stress and compresses at quarter-end when regulatory reporting
constraints bind. Hedge funds supply capacity that is increasing in $\mu_t$ when
the basis exceeds their entry cost but collapses when risk capacity is exhausted
in acute stress. Central-bank swap lines provide explicitly counter-cyclical
capacity by policy design.


% ------------------------------------------------------------------------------
\subsection{Equilibrium}
\label{subsec:equilibrium}
% ------------------------------------------------------------------------------

Each sector $s \in \{FB, USbank, dealer, HF, AM, Corp\}$ takes the funding
spread $\mu_{m,t}$ as given and chooses a signed net USD position
$\mathfrak{q}_{s,m,t}$, with positive values indicating net USD receipt (demand)
and negative values indicating net USD supply. Following \citet{KoijenYogo2019},
sectoral net demand is linear in the spread:
\begin{equation}
  \mathfrak{q}_{s,m,t} = \alpha_{s,m} + \beta_s\,\mu_{m,t} +
  \gamma_s X_{s,m,t} + \varepsilon_{s,m,t},
  \label{eq:net-demand}
\end{equation}
where $\alpha_{s,m}$ is a sector-market fixed effect, $\beta_s$ is sector $s$'s
price sensitivity, $X_{s,m,t}$ is a vector of sector-specific demand shifters,
and $\varepsilon_{s,m,t}$ is a mean-zero residual. Dollar demanders (foreign
banks, corporates, most asset managers) have $\beta_s < 0$; basis traders (hedge
funds acting as suppliers) have $\beta_s > 0$.

Market clearing requires that signed net positions sum to zero in each market:
$\sum_s \mathfrak{q}_{s,m,t} = 0$. Under the maintained hypothesis that
aggregate net demand is downward-sloping, $\mathcal{B} \equiv \sum_s \beta_s <
0$, the market-clearing condition yields a unique equilibrium spread:
\begin{equation}
  \mu_{m,t}^{*} = -\frac{1}{\mathcal{B}} \sum_s \bigl[\alpha_{s,m} +
  \gamma_s X_{s,m,t} + \varepsilon_{s,m,t}\bigr].
  \label{eq:price-eq}
\end{equation}

\begin{theorem}[Equilibrium Existence and Uniqueness]
  \label{thm:equilibrium}
  Under the standard assumptions on preferences, the matching technology, and
  $\mathcal{B} < 0$, a static equilibrium $\mathcal{E}_t =
  (O_t^*, I_t^*, \chi_t^*, p_t^*, \theta_t^*, \mu_t^*)$ exists and is unique
  for any $(U_t, B_t) \in \mathbb{R}_{++}^2$.
\end{theorem}

The proof has three steps. Given $(U_t, B_t)$, market tightness $\theta_t =
U_t / B_t$ determines matching probability $p_t = \eta \theta_t^{-(1-\gamma)}$
and hence quality $\chi_t = \chi(p_t)$. Cost minimization by the representative
agent---who must meet a binding liquidity requirement $\bar{L}_t > 0$ and takes
prices as given---then yields unique demands $(O_t^*, I_t^*)$. Finally,
market clearing in the sectoral demand system pins down the unique $\mu_{m,t}^*$,
because the aggregate net-demand function $F(\mu) = \mathcal{B}\mu + \text{const}$
is strictly decreasing in $\mu$ under $\mathcal{B} < 0$ and crosses zero exactly
once.

The economic intuition is direct. The funding spread $\mu_t^*$ adjusts to clear
the market: when more seekers compete for a given stock of intermediary capacity,
$\theta_t$ rises, matching becomes harder, quality $\chi_t$ falls, and agents
must demand more inside liquidity to meet their payment obligations. Aggregate
net demand rises at any posted spread, and $\mu_t^*$ must increase to restore
clearing. The spread is thus the price that equilibrates the conversion market.

\begin{remark}[Maintained Hypothesis on Aggregate Demand Slope]
  \label{rem:maintained}
  \label{rem:beta-ordering}
  The condition $\mathcal{B} < 0$ is maintained as a hypothesis and tested ex
  post by the sign of the estimated $\hat{\mathcal{B}}$ from the demand system
  in \cref{subsec:demand-system}. If $\hat{\mathcal{B}} \geq 0$, downward-sloping
  aggregate demand is rejected by the data, and the equilibrium price equation
  fails. A negative and significant $\hat{\mathcal{B}}$ confirms the maintained
  hypothesis and provides a structural interpretation of the IV estimates.
  The maintained empirical ordering $|\hat\beta_{FB}| < \hat\beta_{HF}$ follows
  from the model's prediction that basis traders are more price-elastic than
  foreign banks facing inelastic payment obligations.
\end{remark}


% ------------------------------------------------------------------------------
\subsection{Comparative Statics}
\label{subsec:comparative-statics}
% ------------------------------------------------------------------------------

Three propositions govern the response of the equilibrium spread to demand and
capacity shocks. Together they constitute the testable content of the model.

\bigskip
\noindent\textbf{Proposition 1: Demand pressure raises the funding premium.}

\begin{proposition}[Demand and the Funding Spread]
  \label{prop:demand}
  At any interior equilibrium, $\partial \mu_t^* / \partial U_t > 0$: an
  increase in the mass of seekers raises the equilibrium funding spread, holding
  intermediary capacity $B_t$ fixed.
\end{proposition}

The mechanism runs through quality. A rise in $U_t$ tightens the market
($\theta_t \uparrow$), reduces the matching probability ($p_t \downarrow$), and
degrades quality ($\chi_t \downarrow$). Agents facing a binding liquidity
requirement must then demand more inside liquidity at any posted spread, shifting
the aggregate net-demand curve up. The equilibrium spread rises to restore
clearing. Formally, $\partial \mu_t^* / \partial U_t = (\partial \mu_t^* /
\partial \chi_t) \cdot \chi'(p_t) \cdot (\partial p_t / \partial U_t) > 0$,
where each term has a determinate sign: quality deterioration raises the spread
($\partial \mu_t^* / \partial \chi_t < 0$, from the implicit function theorem on
market clearing), quality is increasing in matching probability ($\chi' > 0$),
and matching probability is decreasing in seekers ($\partial p_t / \partial U_t
< 0$). The product is positive.

\paragraph{Testable prediction.} Proposition~1 is the primary empirical target.
The shift-share instrument $Z^{MMF}_{m,t} = \text{Exposure}^{MMF}_{m,t-1}
\times \text{FundingShock}_t$ shifts $U_t$ via the foreign-bank funding channel
while leaving intermediary capacity unaffected, identifying $\partial \mu_t^* /
\partial U_t > 0$. The IV estimate $\hat{\beta} \to -\gamma_{FB}/\mathcal{B}$
under instrument validity, where $\gamma_{FB}$ is the loading of the
foreign-bank sector-specific demand shifter on the aggregate demand shock and
$\mathcal{B}$ is the aggregate slope.

\bigskip
\noindent\textbf{Proposition 2: Capacity tightening amplifies price impact.}

\begin{proposition}[Capacity and the Funding Spread]
  \label{prop:capacity}
  At any interior equilibrium, $\partial \mu_t^* / \partial B_t < 0$: an
  increase in intermediary capacity reduces the equilibrium funding spread,
  holding seeker mass $U_t$ fixed. Furthermore, the price impact of a given
  demand shock is larger when capacity is low:
  \[
    \left.\frac{\partial \mu_t^*}{\partial U_t}\right|_{B_t = B^{low}} \;>\;
    \left.\frac{\partial \mu_t^*}{\partial U_t}\right|_{B_t = B^{high}}.
  \]
\end{proposition}

Higher capacity loosens the market, raises matching probability, improves quality,
and allows the same mass of seekers to be accommodated at a lower spread. The
amplification result---the interaction prediction---follows from the fact that
$|\partial p_t / \partial U_t| = \eta(1-\gamma)\theta_t^{-(2-\gamma)} / B_t$ is
larger when $B_t$ is smaller: the same demand shock produces a larger tightening
of $\theta_t$ when intermediary capacity is already thin.

\paragraph{Testable prediction.} Proposition~2 maps to the interaction
specifications in \cref{subsec:heterogeneity}. The causal effect of $B_t$ cannot
be cleanly identified---dealer constraints tighten precisely in high-stress
states---so the empirical test is stated as effect heterogeneity: the price
impact of instrumented foreign-bank demand is larger when dealers are constrained
(DealerConstraint $= 1$) than when they are not. A positive coefficient on the
$\hat{Q}_{FB,m,t} \times \text{DealerConstraint}_t$ interaction term is
consistent with Proposition~2 but does not establish the causal capacity
channel.

\bigskip
\noindent\textbf{Proposition 3: Convexity of price impact.}

\begin{proposition}[Non-linearity and Threshold Effects]
  \label{prop:nonlinear}
  Under the additional conditions that quality is strictly convex in matching
  probability ($\chi'' > 0$) and that the equilibrium spread responds more than
  proportionally to quality deterioration ($\Phi'' > 0$ where $\Phi$ maps quality
  to the equilibrium spread), the second derivative satisfies
  $\partial^2 \mu_t^* / \partial \theta_t^2 > 0$: the funding spread is a convex
  function of market tightness. Under additional regularity, there exists a
  threshold $\theta^* > 0$ above which the rate of increase $\partial \mu_t^* /
  \partial \theta_t$ escalates sharply.
\end{proposition}

The result is driven by the composition of three effects in the chain
$\theta_t \to p_t \to \chi_t \to \mu_t^*$. Convexity of $p(\theta)$ in
$\theta$ is inherited from the matching function (since
$p''(\theta) = \eta(1-\gamma)(2-\gamma)\theta^{-(3-\gamma)} > 0$). Convexity
of $\chi$ in $p$ (condition $\chi'' > 0$) requires that quality deteriorates
at an accelerating rate as matching becomes difficult. Convexity of $\Phi$ in
$\chi$ (condition $\Phi'' > 0$) requires that the spread responds more than
proportionally to quality deterioration---a condition implied when the market is
thin (small $|\mathcal{B}|$) and quality deterioration reduces the number of
active participants.

\emph{Proposition~3 characterizes the theoretical shape of the price impact
function; testing this result directly would require a threshold specification
beyond the scope of this paper.} The feasibility assessment grades continuous
$\hat{\theta}_t$ measurement at C, and the explorer's report confirms the result
is not testable with the current design. The interaction specifications in
\cref{subsec:heterogeneity} test monotone heterogeneity in price impact across
capacity regimes---a necessary but not sufficient condition for a threshold. We
present Proposition~3 as a theoretical result that motivates our focus on
crisis episodes and quarter-end patterns as high-$\theta_t$ regimes.


% ------------------------------------------------------------------------------
\subsection{Comparative Statics: Quarter-End Amplification}
\label{subsec:comparative-statics-qe}
% ------------------------------------------------------------------------------

\begin{proposition}[Quarter-End Regulatory Amplification]\label{prop:qe}
  Intermediary capacity $B_t$ contracts deterministically at quarter-end
  reporting dates ($Q_t = 1$), lowering $B_t^{FBO}$ and $B_t^{USbank}$ by
  amounts $\omega > 0$ (see \cref{eq:capacity}). By \cref{prop:capacity}, this contraction amplifies
  the price impact of demand shocks:
  $\partial^2 \mu_t^* / \partial U_t \partial Q_t > 0$.
\end{proposition}

The mechanism is the same as Proposition~2 but the supply-side variation is
exogenous to demand: FBO branches reduce their balance sheets at each quarter-end
for regulatory reporting, compressing $B_t^{FBO}$ and $B_t^{USbank}$ by
$\omega > 0$ regardless of the concurrent demand state. The deterministic timing
of this capacity withdrawal distinguishes the quarter-end amplification result
from the endogenous dealer-constraint interaction in \cref{prop:capacity}.

% ------------------------------------------------------------------------------
\subsection{Extensions: Demand System and Structural Interpretation}
\label{subsec:theory-extensions}
\label{subsec:demand-system}
% ------------------------------------------------------------------------------

\begin{remark}[Price-Impact Decomposition and Structural Bridge]\label{rem:price-impact}
The IV estimator $\hat\beta$ converges to $-\gamma_{FB}/\mathcal{B}$ under
instrument validity. This structural bridge connects the reduced-form price
impact to two structural objects: the foreign-bank demand loading $\gamma_{FB}$
and the aggregate demand slope $\mathcal{B}$. Markets with thinner aggregate
price sensitivity ($|\mathcal{B}|$ small) exhibit proportionally larger price
impacts for any given demand shock.
\end{remark}

The static equilibrium in \cref{thm:equilibrium} yields the price equation
$\mu_{m,t}^* = -(1/\mathcal{B}) \sum_s [\alpha_{s,m} + \gamma_s X_{s,m,t} +
\varepsilon_{s,m,t}]$, which connects the IV estimate to structural parameters.
Specifically, the IV estimator $\hat\beta$ converges to $-\gamma_{FB}/\mathcal{B}$
under instrument validity: it is the price impact of a unit shift in foreign-bank
demand, scaled by the aggregate demand slope. This structural interpretation has
two implications. First, the sign of $\hat\beta$ is positive if and only if
$\mathcal{B} < 0$ (downward-sloping aggregate demand) and $\gamma_{FB} > 0$
(foreign-bank demand is increasing in the funding shock). Second, the magnitude
of $\hat\beta$ depends on $\mathcal{B}$: markets with fewer active participants
($|\mathcal{B}|$ small) exhibit larger price impacts for the same demand shock.

In \cref{sec:demand-system}, we embed the equilibrium in a demand system to
recover the aggregate demand elasticity $\mathcal{B} = \sum_s \beta_s$ from
sector-specific price sensitivity estimates. The maintained hypothesis
$\mathcal{B} < 0$ is then testable: a negative and significant
$\hat{\mathcal{B}}$ confirms that the equilibrium pricing equation holds in the
data and provides a structural interpretation of the IV estimates from
\cref{sec:results}.

% ==============================================================================
